% !PRINTED 690
\weberpage{690}
%<<CONT
and if we therefore combine these $w$ associated numbers into one complex, then the number of these complexes is equal to the number $T$ of all non-associated integers in $\Omega$ divisible by $\mathfrak{a}$ whose absolute norm is less than $t$. Accordingly the number of lattice points lying in $S'$ is
\weq{14}{Z_t = w T.}

For the boundary of the region $S$ one obtains from (12) and (13):
\weq{15}{0 \leqq \xi_1 < 1, \ldots, 0 \leqq \xi_{\nu-1} < 1, \xi_\nu < 0,}
and from (5) we obtain
\weq{16}{\mathrm{Lim}\, \frac{T}{t} = \frac{1}{w} \int\!\int \ldots \int dx_1\, dx_2 \ldots dx_n = \frac{1}{w} V,}
in which the $n$-fold integral is to be extended over the region $S$ determined by (15).

\begin{center}
\S.\,190.

\gesperrt{Determination of the Volume.}
\end{center}

The volume $V$ can be determined by the transformation formula for multiple integrals already cited in \S.\,155, according to which, if $x_1, x_2, \ldots, x_n$ are functions of the variables $u_1, u_2, \ldots, u_n$,
\weq{1}{\begin{aligned}
V &= \int\!\int \ldots \int dx_1\, dx_2 \ldots dx_n\\
&= \int\!\int \ldots \int \left( \Sigma \pm \frac{\partial x_1}{\partial u_1} \frac{\partial x_2}{\partial u_2} \cdots \frac{\partial x_n}{\partial u_n} \right) du_1\, du_2 \ldots du_n
\end{aligned}}
holds, where the functional determinant in the integral taken with respect to the variables $u$ is to be taken with its absolute value.

This functional determinant is also equal to the reciprocal value of the determinant
\[ \Sigma \pm \frac{\partial u_1}{\partial x_1} \frac{\partial u_2}{\partial x_2} \cdots \frac{\partial u_n}{\partial x_n}. \]

We first introduce the variables $u$ by a linear substitution, and remark that among the fields conjugate with $\Omega$ there are, according to our assumption, $n - \nu$ conjugate imaginary pairs
%CONT>>
